\section{效果与适用范围}
\label{sec:experiments}

\paragraph{主实验。}
表~\ref{tab:ink-vision-summary} 汇总三个视觉骨干及三种任务规模的实测结果，区分工作台原型与正式实现。正式 \method{} 在 ViT-B/32 八任务上，四个校准 seed 的平均准确率为 $87.059\pm0.021\%$，其中误差为样本标准差；比已复现的 ESM 高约 $2.87$ 个百分点。工作台原型九格均已完成，B/32 八任务为 $86.60\%$；正式实现尚未补齐其余八格，因此分别列行。逐任务归一化均使用对应环境的实测专家，正式环境专家宏平均为 $89.96\%$。主表缺测格不参与方法排名。

% Generated by motivation_method/update_results.py from data/verified_results.json.
\begin{table}[!htbp]
\centering
\caption{CLIP 实测结果（\%）。单元格为任务宏平均准确率，下标为逐任务相对本环境专家的归一化均值。上半表为工作台；正式 CM 单独列出。--- 为缺测。}
\label{tab:ink-vision-summary}
\setlength{\tabcolsep}{2pt}
\renewcommand{\arraystretch}{1.12}
\resizebox{\textwidth}{!}{\begin{tabular}{@{}lccc|ccc|ccc@{}}
\toprule
\multirow{2}{*}{Method} & \multicolumn{3}{c}{ViT-B/32} & \multicolumn{3}{c}{ViT-B/16} & \multicolumn{3}{c}{ViT-L/14} \\
\cmidrule(lr){2-4}\cmidrule(lr){5-7}\cmidrule(l){8-10}
 & 8 tasks & 14 tasks & 20 tasks & 8 tasks & 14 tasks & 20 tasks & 8 tasks & 14 tasks & 20 tasks \\
\midrule
Individual FT & $89.96_{(100.00)}$ & $88.88_{(100.00)}$ & $89.76_{(100.00)}$ & $91.58_{(100.00)}$ & $90.62_{(100.00)}$ & $91.52_{(100.00)}$ & $93.11_{(100.00)}$ & $91.96_{(100.00)}$ & $92.87_{(100.00)}$ \\
\midrule
Zero-shot & $46.69_{(53.63)}$ & $56.05_{(63.70)}$ & $54.67_{(61.95)}$ & $53.10_{(58.89)}$ & $59.62_{(65.85)}$ & $57.91_{(63.73)}$ & $61.28_{(66.52)}$ & $66.12_{(71.83)}$ & $63.36_{(68.56)}$ \\
Weight Avg. & $66.42_{(73.83)}$ & $62.75_{(70.77)}$ & $59.22_{(66.72)}$ & $72.36_{(78.72)}$ & $67.29_{(73.87)}$ & $62.74_{(68.70)}$ & $80.33_{(86.06)}$ & $75.33_{(81.48)}$ & $70.07_{(75.44)}$ \\
Task Arith. & $66.49_{(73.76)}$ & $62.71_{(70.75)}$ & $58.98_{(66.49)}$ & $72.30_{(78.58)}$ & $66.99_{(73.48)}$ & $62.46_{(68.36)}$ & $80.37_{(86.09)}$ & $75.29_{(81.45)}$ & $69.57_{(74.83)}$ \\
TIES & $69.41_{(76.94)}$ & $64.53_{(72.59)}$ & $60.81_{(68.31)}$ & $74.49_{(81.01)}$ & $68.89_{(75.47)}$ & $64.84_{(70.82)}$ & $82.15_{(88.02)}$ & $76.17_{(82.41)}$ & $72.32_{(77.74)}$ \\
Consensus TA & $70.36_{(77.86)}$ & $67.70_{(75.97)}$ & $63.27_{(70.84)}$ & $76.08_{(82.56)}$ & $72.33_{(79.26)}$ & $67.77_{(73.97)}$ & $82.40_{(88.21)}$ & $78.43_{(84.79)}$ & $75.66_{(81.24)}$ \\
Iso-C & $81.33_{(90.21)}$ & $75.93_{(85.18)}$ & $71.18_{(79.36)}$ & $86.26_{(93.87)}$ & $80.79_{(88.65)}$ & $75.77_{(82.53)}$ & $90.43_{(97.05)}$ & $85.15_{(92.16)}$ & $83.50_{(\text{---})}$ \\
Iso-CTS & $81.04_{(89.86)}$ & $77.82_{(87.31)}$ & $73.79_{(82.16)}$ & $87.47_{(95.27)}$ & $82.71_{(90.84)}$ & $78.78_{(85.87)}$ & $91.29_{(97.99)}$ & $86.90_{(94.16)}$ & $86.21_{(\text{---})}$ \\
TSV-M & $81.35_{(90.08)}$ & $76.69_{(85.93)}$ & $73.89_{(82.23)}$ & $85.01_{(92.45)}$ & $80.89_{(88.85)}$ & $77.35_{(84.31)}$ & $88.57_{(94.88)}$ & $85.50_{(92.72)}$ & $84.26_{(\text{---})}$ \\
RegMean & $81.27_{(90.00)}$ & $76.14_{(85.27)}$ & $72.65_{(80.99)}$ & $84.11_{(91.50)}$ & $79.27_{(86.94)}$ & $75.96_{(82.80)}$ & $87.59_{(93.79)}$ & $83.80_{(90.57)}$ & $81.66_{(87.61)}$ \\
WUDI & $82.55_{(91.37)}$ & $75.41_{(84.53)}$ & $66.30_{(74.19)}$ & $85.54_{(93.00)}$ & $78.83_{(86.43)}$ & $70.25_{(76.61)}$ & $90.11_{(96.60)}$ & $85.44_{(92.69)}$ & $80.76_{(86.77)}$ \\
LiNeS (Iso-C) & $83.37_{(92.58)}$ & $78.09_{(87.64)}$ & $73.82_{(82.22)}$ & $87.27_{(95.09)}$ & $82.21_{(90.32)}$ & $77.60_{(84.54)}$ & $90.88_{(97.54)}$ & $86.17_{(93.36)}$ & $84.56_{(90.76)}$ \\
CoM & $83.38_{(92.30)}$ & $78.46_{(87.76)}$ & $75.23_{(83.76)}$ & $85.56_{(93.13)}$ & $80.93_{(88.80)}$ & $77.27_{(84.20)}$ & $88.84_{(95.17)}$ & $85.29_{(92.26)}$ & $83.68_{(89.80)}$ \\
ESM & $84.19_{(93.48)}$ & $80.97_{(90.99)}$ & $78.95_{(87.88)}$ & $88.06_{(96.05)}$ & $84.49_{(93.14)}$ & $81.94_{(89.52)}$ & $91.01_{(97.69)}$ & $87.57_{(95.09)}$ & --- \\
SVC (Iso-CTS) & $81.61_{(90.79)}$ & $77.79_{(87.31)}$ & $73.91_{(82.45)}$ & $87.37_{(95.28)}$ & $82.68_{(90.82)}$ & $78.67_{(85.82)}$ & $90.64_{(97.31)}$ & $86.57_{(93.73)}$ & $85.37_{(91.69)}$ \\
SVC (ESM) & $82.38_{(91.58)}$ & $79.03_{(88.83)}$ & $76.62_{(85.49)}$ & $86.51_{(94.15)}$ & $82.47_{(90.67)}$ & $79.74_{(86.95)}$ & $89.79_{(96.29)}$ & $86.76_{(94.10)}$ & --- \\
\rowcolor{black!8}
CM（原型） & $86.60_{(96.13)}$ & $83.93_{(94.17)}$ & $83.42_{(92.80)}$ & $88.99_{(97.04)}$ & $86.45_{(95.15)}$ & $85.83_{(93.68)}$ & $91.52_{(98.24)}$ & $88.75_{(96.29)}$ & $88.96_{(95.64)}$ \\
\midrule
\rowcolor{black!8}
CM（正式，4 seeds） & $87.06_{(96.64)}$ & --- & --- & --- & --- & --- & --- & --- & --- \\
\bottomrule
\end{tabular}}
\par\smallskip
{\footnotesize 工作台 CM 九格均为单种子，按无标签留出 KL 选轮；可调基线按有标签 validation 选参。L/14 二十任务的 Iso-C、Iso-CTS、TSV-M 仅恢复了日志中的两位小数宏平均，故归一化下标留空。正式 CM 为 $87.059\pm0.021$，使用正式运行中的专家作归一化分母；其余行使用工作台闸门专家。两版不拼成一行。}
\end{table}


表~\ref{tab:ink-t5-results} 给出 Flan-T5-base 八个任务专家一次性合并后的逐任务实测结果，采用 FusionBench 风格的 GLUE validation 生成评测；它不是持续合并实验。原版 \method{} 的五种子宏平均为 $84.92\pm0.05$；目标插值扩展 $0.3$ 为四种子 $85.64\pm0.04$，扩展 $0.6$ 为三种子 $85.77\pm0.08$。移植 FeatCal 的五种子结果为 $84.84\pm0.22$，统一报告样本标准差。原版与扩展分别报告，不将不同设定上的最高分拼为同一方法版本。

% Generated by motivation_method/update_results.py from data/verified_results.json.
\begin{table}[!htbp]
\centering
\caption{Flan-T5-base 八专家一次性合并后的 GLUE validation 实测结果（FusionBench 口径）。前七任务为生成答案 exact match，STSB 为 Spearman；末列为八任务宏平均。}
\label{tab:ink-t5-results}
\setlength{\tabcolsep}{3pt}
\renewcommand{\arraystretch}{1.14}
\resizebox{\textwidth}{!}{\begin{tabular}{@{}lrrrrrrrr|c@{}}
\toprule
Method & CoLA & MNLI & MRPC & QNLI & QQP & RTE & SST2 & STSB & 平均 \\
\midrule
Individual FT & 74.98 & --- & --- & --- & --- & --- & --- & --- & --- \\
Zero-shot & 69.13 & 56.45 & 76.23 & 88.45 & --- & --- & --- & --- & --- \\
\midrule
Weight Avg. & 69.13 & 62.65 & 79.41 & 89.80 & 83.86 & 81.23 & 91.74 & 73.21 & $78.88$ \\
Task Arith. & 70.47 & 57.77 & 78.43 & 90.21 & 83.62 & 80.51 & 92.32 & 77.82 & $78.89$ \\
TIES & 70.09 & 66.25 & 78.43 & 90.12 & 83.62 & 81.59 & 91.74 & 78.24 & $80.01$ \\
Iso-C 1.3 & 69.51 & 67.27 & 78.43 & 90.17 & 84.38 & 81.23 & 92.20 & 73.96 & $79.64$ \\
TSV-M & 70.76 & 79.74 & 78.68 & 90.68 & 84.27 & 83.39 & 93.12 & 81.44 & $82.76$ \\
RegMean 0.9 & 69.80 & 83.32 & 79.66 & 90.90 & 84.51 & 83.39 & 93.12 & 85.71 & $83.80$ \\
TA + FeatCal（移植） & 72.06 & 82.50 & 83.82 & 90.94 & 84.80 & 84.19 & 93.00 & 87.38 & $84.84\pm 0.22$ \\
\rowcolor{black!8}
\method{}（原版） & 71.91 & 83.38 & 84.26 & 91.22 & 84.79 & 83.18 & 93.21 & 87.38 & $84.92\pm 0.05$ \\
\rowcolor{black!8}
\method{}（目标插值 0.3） & 74.98 & 83.25 & 85.54 & 90.94 & 84.87 & 84.39 & 93.26 & 87.86 & $85.64\pm 0.04$ \\
\rowcolor{black!8}
\method{}（目标插值 0.6） & 75.07 & 82.95 & 85.70 & 90.79 & 84.80 & 85.32 & 93.31 & 88.18 & $85.77\pm 0.08$ \\
\bottomrule
\end{tabular}}
\par\smallskip
{\footnotesize 多种子逐任务列为种子均值；末列误差为各次八任务均值的样本标准差。FeatCal 与原版 CM 各 5 seeds，目标插值 0.3 为 4 seeds，0.6 为 3 seeds；其他合并行为单种子。Zero-shot 仅测四任务，独立专家仅有 CoLA 对角结果，均不补八任务平均；未完成 T5 适配的视觉基线不列入本表。}
\end{table}


\paragraph{跨任务规模与模型设定。}
在独立的 CLIP 工作台实现中，三个骨干与 $8/14/20$ 任务构成的九格均较固定 Iso-C 起点提高约 $1.38$--$16.66$ 个百分点（表~\ref{tab:clip-scaling}）。任务规模变化同时改变任务集合，不能将宏平均变化只归因于任务数。T5-large 原版为 $88.99$，目标插值 $0.3$ 为 $87.97$，均为单种子；扩展在 base 上的收益没有迁移到 large（表~\ref{tab:t5-main}）。T5 使用教师答案上的 teacher-forced 分布，LLM 使用教师 top-128 logits；两者保留“当前统计、逐层求解、按 $\mathcal L_{\mathcal H}$ 选步、接受后重测”的共同流程，具体教师近似和搜索粒度按设定适配。LLM 十组起点对照中有九组 GSM8K 提高（表~\ref{tab:llm-summary}及完整附录表）：Llama 的 RegMean 起点下降，TSV-M 起点的 IFEval 也下降；Llama 的三域宏平均最高仍是 TSV-M 起点自身 $39.97$，接受 CM 后为 $39.26$。Gemma 的 Localize-and-Stitch 为 $51.33$，高于所列后处理结果。因此，我们主张的是\emph{同一迭代框架在多种设定中可用}，而不是所有任务均提升、所有基线均被超越，或完全不需设定适配。

\paragraph{起点与组件性质。}
初始化适配作为附加性质检验：ViT-B/32 的 Average 与 RegMean 诊断准确率差距由 $22.36$ 缩到 $0.15$ 个百分点（图~\ref{fig:motivation}a），但未拉平对照仍保留明显差距（附录~\ref{app:protocol}）。把 \method{} 叠在其它合并输出上的工作台结果见附录~\ref{app:cross-setting}，与图~\ref{fig:motivation}a 的诊断实验不是同一协议。归一化的同机直接消融为 $87.0451\%$ 对 $86.6652\%$，其几何解释及优势条件见附录~\ref{app:normalization}。这些结果分别检验起点适配与组件作用，并与图~\ref{fig:motivation}b/c 的迭代和度量诊断互补。
